\documentclass[11pt]{article}
\usepackage[letterpaper, margin=0.85in, top=0.75in, bottom=0.8in]{geometry}
\usepackage{libertinus}
\usepackage[scale=0.82]{sourcecodepro}
\usepackage{xcolor, fvextra, booktabs, enumitem, titlesec, microtype, tabularx}
\usepackage[most]{tcolorbox}
\usepackage[hidelinks]{hyperref}

% J.T.S. palette
\definecolor{ink}{HTML}{1E2B24}
\definecolor{muted}{HTML}{56645B}
\definecolor{accent}{HTML}{2B5E8C}
\definecolor{accentsoft}{HTML}{DCE6EF}
\definecolor{moss}{HTML}{4F6E4C}
\definecolor{sand}{HTML}{CDBF9E}
\definecolor{panel}{HTML}{F7F9F5}
\definecolor{rulec}{HTML}{D2DACD}
\color{ink}

\newcommand\Kw[1]{\textcolor{accent}{\textbf{#1}}}
\newcommand\Cm[1]{\textcolor{muted}{#1}}
\newcommand\Rl[1]{\textcolor{accent}{\textbf{#1}}}
\newcommand\Ms[1]{\textcolor{moss}{\textbf{#1}}}
\newcommand\Dm[1]{\textcolor{muted}{#1}}
\definecolor{mosssoft}{HTML}{E1EADC}
\newcommand\Qd[1]{\colorbox{mosssoft}{\textcolor{moss}{\textbf{#1}}}}
\newtcolorbox{termbox}{breakable, enhanced, colback=panel, colframe=rulec, boxrule=0.4pt,
  borderline west={2.2pt}{0pt}{sand}, arc=2pt, left=7pt, right=5pt, top=3pt, bottom=3pt,
  before skip=4pt, after skip=8pt}

\fvset{fontsize=\small, breaklines, commandchars=\\\{\}}
\newtcolorbox{thmbox}[1]{enhanced, colback=panel, colframe=rulec, boxrule=0.4pt,
  borderline west={2.4pt}{0pt}{accent}, arc=2pt, left=8pt, right=6pt, top=5pt, bottom=5pt,
  title={#1}, coltitle=accent, fonttitle=\sffamily\bfseries\small, attach title to upper={\par\smallskip},
  before skip=6pt, after skip=8pt}

\titleformat{\section}{\sffamily\large\bfseries\color{ink}}{}{0pt}{}
\titlespacing*{\section}{0pt}{14pt}{4pt}
\setlength\parindent{0pt}
\setlength\parskip{5pt}
\setlength\emergencystretch{3em}
\pagestyle{plain}

\makeatletter
% typewriter identifiers are not hyphenated or broken at their hyphens
\AddToHook{selectfont}{\edef\jts@f{\f@family}\edef\jts@t{\ttdefault}%
  \ifx\jts@f\jts@t\hyphenchar\font=-1\relax\fi}
\makeatother

\begin{document}

{\sffamily\small\bfseries\color{moss}PVS\_CAD \quad·\quad PVS 8.1 + NASALib \quad·\quad OCTOBER 2026}\par\vspace{6pt}
{\fontsize{24}{29}\selectfont\bfseries A verified CAD algorithm and\\ decision procedure in PVS}\par\vspace{8pt}
{\large\color{muted}An overview of the pvs\_cad library and its proof strategies \texttt{(cad)} and \texttt{(cad-qe)}}\par\vspace{8pt}
{\large J.~Tanner Slagel{\renewcommand\thefootnote{*}\footnote{This report was written with Claude Code using Claude Opus 5.5. The pvs\_cad library was developed the same way, with Claude Fable 5.1 and then Opus 5 and Opus 5.5; every proof in it is checked by PVS. Claude Code: \url{https://claude.com/claude-code}.}}}\par\vspace{2pt}
{\small\color{muted}Code and proofs: \url{https://github.com/j-tanner-slagel/pvs_cad}}\par\vspace{10pt}

\textbf{In short.} pvs\_cad is a PVS library, with the proof strategies \texttt{(cad)} and \texttt{(cad-qe)}, that
decides the first-order theory of the real numbers and eliminates quantifiers: formulas built from polynomial equations
and inequalities with rational coefficients, the connectives AND, OR, NOT, IMPLIES and IFF, and quantifiers over real
variables. Two executable PVS functions compute a cylindrical algebraic decomposition (CAD). For the cells built from
Collins's projection operator the library proves, for every family of polynomials in every number of variables and with
no run-time certificate, that the cells of $\mathbb{R}^n$ cover it and are disjoint, that the cells of every level are
connected and first-order definable and stack in cylinders over continuous, strictly ordered root functions with
definable graphs, and that each cell of $\mathbb{R}^n$ carries one sign vector of the input polynomials
(\texttt{col\_cells}); \texttt{col\_found} returns this CAD as data, at least one record with an exact sample point for
every nonempty cell of $\mathbb{R}^n$ and none for any other, and those records decide every sentence over the
polynomials in the same variables, in the same order. The n-level engine, built first, computes such a
CAD in two or more variables whenever its run-time certificate holds, which some value of its search-effort parameter
$u$ always achieves (\texttt{cad\_verified}). A set is definable by a first-order formula with rational coefficients
exactly when a quantifier-free one defines it (\texttt{fod\_qfd}, Tarski--Seidenberg over $\mathbb{Q}$), and every cell
of both CADs, at every level, is quantifier-free definable (\texttt{cad\_sa}). In a theory that imports
\texttt{pvs\_cad}, \texttt{(cad)} decides a closed prenex formula over \texttt{real} whose atoms compare polynomials in its
bound variables by \texttt{decide8}, the Collins run (after a witness search), and any other first-order formula over the
reals, one with other real terms included, by \texttt{decide\_g}; without \texttt{cad\_decide8} it uses \texttt{decide5},
built on the n-level engine.
All three are proved \emph{sound} (an answer is always the truth) and \emph{complete} (they always answer).
\texttt{(cad-qe)} replaces a formula with free variables by a quantifier-free one that holds at exactly the same points.
Inside a proof, the strategies evaluate the procedure and turn its answer into a PVS proof that the kernel checks. The
whole library, 355 theories and 5,183 formulas, is proved.

\textbf{Using it.} A theory needs one line, \texttt{IMPORTING pvs\_cad}. Then, in a proof:
\begin{center}\small
\begin{tabularx}{\textwidth}{@{}>{\raggedright\arraybackslash}p{0.29\textwidth} X@{}}
\toprule
\mbox{\texttt{(cad)}} & decide the goal, a first-order sentence over the reals: a TRUE goal is proved; otherwise it stays,
labelled \texttt{cad}, with a message saying it is FALSE\\
\mbox{\texttt{(cad -1)}}, \mbox{\texttt{(cad 2)}}, \mbox{\texttt{(cad (-1 3))}} & the same on a chosen hypothesis or goal (a FALSE hypothesis
closes the goal), or on several formulas together\\
\mbox{\texttt{(cad *)}}, \mbox{\texttt{(cad +)}}, \mbox{\texttt{(cad -)}} & decide the sequent as a whole (the hypotheses imply the goals), or
the goals or the hypotheses alone; other real terms, such as $f(c)$ or $\sqrt{c}$, become unknowns\\
\mbox{\texttt{(cad-qe)}}, \mbox{\texttt{(cad-qe -1)}} & replace a goal or hypothesis with free real terms by an equivalent
quantifier-free formula\\
\mbox{\texttt{(cad :cad-only? t)}}, \mbox{\texttt{(cad-direct)}} & answer only through the verified CAD engine; decide without the
witness search\\
\bottomrule
\end{tabularx}
\end{center}
Formulas are taken as written: grouped binders, quantifiers inside NOT, AND, OR, IMPLIES and IFF, PVS's synonyms for the
connectives, and binders over \texttt{posreal}, \texttt{nnreal}, \texttt{negreal}, \texttt{npreal} and \texttt{nzreal}.
\texttt{cad/cad\_showcase.pvs} has 43 commented examples, each with its proof and its time, and
\texttt{cad/cad\_forms\_ex.pvs} 64 regression lemmas.

\section{What it looks like}
Real PVS sessions, recorded on 1 October 2026 in \texttt{cad\_showcase}, which imports \texttt{pvs\_cad}, with the
prover's commentary and the simplifier's internal messages cut. The quadratic formula as an existence theorem, written
with a grouped binder (\texttt{(cad)} proves it from its nested form; the timing table below gives the whole proof of
that nested form, \texttt{d\_quad\_suf}):
\begin{termbox}
\begin{Verbatim}[commandchars=\\«», fontsize=\footnotesize]
\Dm«quadratic_formula :»
\Dm«  |-------»
\Dm«{1}»   FORALL (a, b, c: real):
        EXISTS (x: real):
          a /= 0 AND b ^ 2 >= 4 * a * c IMPLIES a * x ^ 2 + b * x + c = 0
\Rl«Rule? (cad)»
\Ms«Deciding 1 by cylindrical decision, witness first,»
\Qd«Q.E.D.»
\Dm«Real time = 5.297 secs.»
\end{Verbatim}
\end{termbox}

Continuity of $x^2$ at 0, the epsilon--delta way: three alternating quantifiers, binders over \texttt{posreal} and an
implication inside, taken as written:
\begin{termbox}
\begin{Verbatim}[commandchars=\\«», fontsize=\footnotesize]
\Dm«sq_continuous :»
\Dm«  |-------»
\Dm«{1}»   FORALL (e: posreal):
        EXISTS (d: posreal):
          FORALL (x: real): -d < x AND x < d IMPLIES x ^ 2 < e
\Rl«Rule? (cad)»
\Ms«Deciding 1 by cylindrical decision, witness first,»
\Qd«Q.E.D.»
\Dm«Real time = 1.81 secs.»
\end{Verbatim}
\end{termbox}

\texttt{(cad *)} works on a whole sequent. After \texttt{(skeep)} the facts sit in the sequent over skolem constants;
\texttt{f(x)} and \texttt{length(l)} become unknown reals, and that \texttt{length(l)} is nonnegative comes from its
type. The strategy says which formulas it uses, which terms become unknowns, and which sentence it decides:
\begin{termbox}
\begin{Verbatim}[commandchars=\\«», fontsize=\footnotesize]
\Dm«unknown_terms :»
\Dm«  |-------»
\Dm«{1}»   FORALL (f: [real -> real], l: list[real], x: real):
        f(x) > x ^ 2 AND length(l) = 3 IMPLIES f(x) + length(l) > 2
\Rl«Rule? (skeep)»
\Dm«{-1}»  f(x) > x ^ 2
\Dm«{-2}»  length(l) = 3
\Dm«  |-------»
\Dm«{1}»   f(x) + length(l) > 2
\Rl«Rule? (cad *)»
\Ms«cad*: using (-1 -2 1), skipping nothing; unknown reals: cadv_1 = f(x), cadv_2 = x, cadv_3 = length(l)»
\Ms«cad*: deciding FORALL (cadv_1: real): FORALL (cadv_2: real): FORALL (cadv_3: real): (((cadv_3 >= 0)) IMPLIES (((cadv_1 > (cadv_2 ^ 2)) AND (cadv_3 = (3))) IMPLIES (((cadv_1 + cadv_3) > (2)))))»
\Ms«Deciding * by cylindrical decision, witness first,»
\Qd«Q.E.D.»
\Dm«Real time = 2.292 secs.»
\end{Verbatim}
\end{termbox}

Quantifier elimination inside a proof: \texttt{(cad-qe -1)} prints the quantifier-free answer and replaces the
hypothesis by it, with a proof that the two are equivalent; \texttt{(cad *)} finishes:
\begin{termbox}
\begin{Verbatim}[commandchars=\\«», fontsize=\footnotesize]
\Dm«qe_then_cad :»
\Dm«  |-------»
\Dm«{1}»   FORALL (a, b: real):
        (EXISTS (t: real): t ^ 2 + a * t + b = 0) AND b > 0 IMPLIES
         a ^ 2 >= 4 * b
\Rl«Rule? (skeep)»
\Dm«{-1}»  EXISTS (t: real): t ^ 2 + a * t + b = 0
\Dm«{-2}»  b > 0
\Dm«  |-------»
\Dm«{1}»   a ^ 2 >= 4 * b
\Rl«Rule? (cad-qe -1)»
\Ms«cad-qe: ((b < 0) OR ((b = 0) OR ((b > 0) AND ((4*b + (-1)*a^2 = 0) OR (4*b + (-1)*a^2 < 0)))))»
\Ms«Eliminating the quantifiers of -1,»
\Dm«{-1}»  ((b < 0) OR
        ((b = 0) OR
          ((b > 0) AND
            ((4 * b + (-1) * a ^ 2 = 0) OR (4 * b + (-1) * a ^ 2 < 0)))))
\Dm«{-2}»  b > 0
\Dm«  |-------»
\Dm«{1}»   a ^ 2 >= 4 * b
\Rl«Rule? (cad *)»
\Ms«cad*: using (-1 -2 1), skipping nothing; unknown reals: cadv_1 = b, cadv_2 = a»
\Ms«cad*: deciding FORALL (cadv_1: real): FORALL (cadv_2: real): ((((cadv_1 < (0)) OR ((cadv_1 = (0)) OR ((cadv_1 > (0)) AND (((((4) * cadv_1) + ((-1) * (cadv_2 ^ 2))) = (0)) OR ((((4) * cadv_1) + ((-1) * (cadv_2 ^ 2))) < (0)))))) AND (cadv_1 > (0))) IMPLIES (((cadv_2 ^ 2) >= ((4) * cadv_1))))»
\Ms«Deciding * by cylindrical decision, witness first,»
\Qd«Q.E.D.»
\Dm«Real time = 4.596 secs.»
\end{Verbatim}
\end{termbox}

\section{How a proof by \texttt{(cad)} works}
The strategy follows the reflection template of Narkawicz, Muñoz and Dutle's univariate Sturm and Tarski strategies in
NASALib~\cite{nasalib,nmd}. \texttt{(cad)} reads the goal and computes its encoding as data \texttt{(os, F, phi)}: the
quantifier prefix, the polynomials' coefficient lists and the Boolean skeleton (see ``The main theorems''). It then runs
the decision in PVS's ground evaluator (\texttt{decide8}, the Collins run, in a theory that imports
\texttt{cad\_decide8}, as \texttt{pvs\_cad} does; \texttt{decide5} in one that imports \texttt{cad\_decide5} but not
\texttt{cad\_decide8}) and
instantiates its correctness theorem (\texttt{decide8\_correct} or \texttt{decide5\_correct}) with that data. Unfolding
\texttt{fsem}, the semantics, turns the result back into the original formula: each quantifier is matched, and one
reflection lemma per polynomial (\texttt{peval\_pnorm}) shows that the coefficient list evaluates to the polynomial as
written. A FALSE sentence is reported and left in place, and a false hypothesis closes the goal. Before deciding a
prenex sentence, \texttt{(cad)} looks for a witness: a model search proves a goal whose quantifiers are all existential
(or refutes a hypothesis whose quantifiers are all universal) by instantiating a point, and with three or more
quantifiers, when the decision is slow, one variable (existential in a goal, universal in a hypothesis being refuted)
is fixed to a small constant and the smaller sentence decided. Grouped binders such as \texttt{FORALL (x, y: real)} are decided through the nested form, and a goal
\texttt{NOT A} as the hypothesis \texttt{A}. In a theory that imports \texttt{pvs\_cad}, a formula of any other shape
(quantifiers inside connectives; binders over \texttt{posreal}, \texttt{nnreal}, \texttt{negreal}, \texttt{npreal} or
\texttt{nzreal}) is encoded as a term \texttt{g} of a datatype that mirrors PVS's connectives; \texttt{(cad)} evaluates \texttt{decide\_g(g)}, instantiates \texttt{decide\_g\_ok}, unfolds the semantics
\texttt{gsem}, which then has the formula's own shape, and closes the goal connective by connective. \texttt{(cad-qe)}
does the same for a formula with free real terms (skolem constants, say): it evaluates \texttt{qe8} (\texttt{qe} in a
theory that does not import \texttt{qe8}), or \texttt{qe\_g} for other shapes, and replaces the formula by the
quantifier-free answer, proved equivalent from \texttt{qe8\_correct}, \texttt{qe\_correct} or \texttt{qe\_g\_ok}.
\texttt{(cad *)} collects the polynomial hypotheses and goals of a whole sequent, replaces non-polynomial real terms
such as \texttt{sqrt(x)} or \texttt{f(x)} by universally quantified reals, and decides the implication, which enters
the proof by \texttt{case}, so both branches are checked by PVS.

The trusted base is the PVS kernel plus the ground evaluator that runs the procedure (\texttt{decide8},
\texttt{decide\_g}, \texttt{qe8} or \texttt{qe\_g} in a theory that imports \texttt{pvs\_cad}; \texttt{decide5},
\texttt{decide7} or \texttt{qe} in one that imports only \texttt{cad\_decide5}, \texttt{cad\_decide7} and
\texttt{qe\_all}), the same base as NASALib's Sturm and Tarski strategies. It also includes NASALib's saved proofs of the
theorems the library imports, which were not replayed here. Beyond that, one trusts the specification: that the
semantics (\texttt{fsem}, \texttt{gsem}, \texttt{qfsem}) and the theorem statements (last section) say what is intended.
The strategy code (3,371 lines of Lisp) is not trusted: it only chooses proof steps.

\section{What it decides in practice}
CAD is doubly exponential in the worst case, so completeness means ``always answers, given time''. Measured times for
complete proofs on one M-series Mac: with \texttt{decide5}, in theories that import \texttt{cad\_decide5}, on
28 September 2026 (PVS's own timing); with \texttt{IMPORTING pvs\_cad} (\texttt{decide8}, the Collins run) on
1 October 2026 (wall clock of the proof command, through \texttt{pvs-cli}):
\begin{center}\small
\begin{tabularx}{\textwidth}{@{}X l r r@{}}
\toprule
Sentence & Shape & \texttt{decide5} & \texttt{pvs\_cad}\\
\midrule
Wilkinson's $(x-1)\cdots(x-20)$, expanded, is positive for $x<1$ (coefficients to $2.4\times10^{18}$) & $\forall$, degree 20 & 2.9\,s & 3.2\,s\\
Chebyshev $T_{10}$ stays in $[-1,1]$ on $[-1,1]$ & $\forall$, degree 10 & 2.0\,s & 2.7\,s\\
Motzkin's polynomial is nonnegative (not a sum of squares~\cite{motzkin}) & $\forall\forall$, degree 6 & 1.1\,s & 1.7\,s\\
Every depressed cubic $x^3+px+q$ has a real root & $\forall\forall\exists$ & 0.8\,s & 1.5\,s\\
Quadratic formula: $a\neq0 \wedge b^2\ge 4ac \Rightarrow \exists x\colon ax^2+bx+c=0$ & $\forall\forall\forall\exists$ & 2.1\,s & 3.2\,s\\
$\forall x\,\exists y\,\forall z\colon z^4+yz^2+y\ge x$ & $\forall\exists\forall$ & 22.9\,s & 8.3\,s\\
No real root of $z^2+bz+c$ implies $b^2<4c$, by \texttt{(cad *)} on a sequent & quantified hypothesis & 1.3\,s & 2.0\,s\\
\bottomrule
\end{tabularx}
\end{center}
With \texttt{decide5} the limits show up with three or four variables at degree 3 to 4: AM--GM for three or four
variables and Cauchy--Schwarz in the plane did not finish in two minutes (28 September 2026). With
\texttt{pvs\_cad} AM--GM for three variables takes 3.4\,s and Cauchy--Schwarz in the plane 2.2\,s, while AM--GM for
four variables still does not finish in two minutes (1 October 2026). On 100 generated sentences with 1 to 6 quantifiers, \texttt{(cad)} with \texttt{decide5}
proved 98 within 300\,s, and no answer disagreed with Z3 (27 September 2026). On 68 \texttt{(cad)} and \texttt{(cad-qe)}
proofs from the library's example theories, proofs with the Collins imports (\texttt{decide8}, \texttt{qe8}) took
0.6--0.8\,s longer on average than with \texttt{decide5} and \texttt{qe} (0.77\,s in a first run, 0.59\,s in a re-run with the
strategy fixes of this release), partly because of the larger imported context;
on some problems with several quantifiers the Collins decision was faster by more than an order of magnitude (23 to
more than 450 times) and on others slower (1 October 2026; \texttt{docs/collins\_cad.pdf}).

\section{How it was built}
The library was built in stages from 9 September to 2 October 2026 (298
commits). Most stages also shipped a proof strategy of their own.
\begin{enumerate}[leftmargin=1.3em, itemsep=2pt, topsep=2pt]
\item \textbf{Multivariate polynomials.} A normal form for multivariate polynomials over the rationals, proved unique, and
reflection lemmas connecting it to PVS real expressions (strategies \texttt{mpoly-eq}, \texttt{mpoly-simp}).
\item \textbf{Real algebraic numbers.} Root isolation, ordering and signs at algebraic numbers, on top of NASALib's
Sturm and Tarski libraries (\texttt{alg-roots}).
\item \textbf{Branching quantifier elimination.} One-variable quantifier elimination with parametric sign conditions
(\texttt{qe-exists}, \texttt{poly-sign}).
\item \textbf{Resultants and subresultants}: determinants, the resultant, and the principal subresultant coefficients
with their specialization property. Two projection operators built from them are defined but not proved delineating:
a McCallum-shaped one (\texttt{cad\_proj}), from which the decisions of item 6 start their read closures (their
soundness rests on run-time checks instead), and a Collins-like one without reducta (\texttt{cad\_projc}). Collins's
own operator, with reducta (the N-reducta projection), is proved delineating in item 10.
\item \textbf{A first cylindrical decision}, \texttt{decide}, proved correct for every sentence
(\texttt{decide\_correct}). It was complete but exponential in the number of polynomials: a circle and a line did not
finish in 15 minutes.
\item \textbf{The fast cylindrical decision}, which \texttt{decide5} uses for two or more quantifiers; \texttt{(cad)}
runs \texttt{decide5} in a theory that does not import \texttt{cad\_decide8}. Over each cell of the lower variables, it
lifts to the next variable by walking exact rational separators between that variable's roots, and reads the sign
vector at each separator and root through Tarski queries at the cell's sample point, which may be algebraic. Soundness
never assumes what it cannot check: the signs each lifting step reads are certified constant on each sector at run
time, which is proved to fix the set of sign vectors in every fibre over the sector, and a failed check sets
\texttt{ok} to FALSE (\texttt{decw\_correct} for two quantifiers, \texttt{decn\_correct} for any number).
\item \textbf{Completeness.} A review showed that the only way \texttt{ok} could be FALSE was a fixed search limit
running out. Some searches were made total; the others share one fuel parameter $u$, and every large enough $u$ makes
every check pass: each search stops at its first success, so a larger $u$ changes nothing once $u$ suffices, and the
closures that add sign reads draw them from a finite set built level by level from \texttt{F}, so they reach a fixpoint
(\texttt{decn\_complete}). \texttt{decn\_u}, which \texttt{decide5} calls, raises
$u$ until the checks pass; it terminates on a non-computable measure, the distance to a bound beyond which every $u$
suffices (\texttt{ubound}, chosen by \texttt{epsilon}), which keeps it an ordinary executable PVS function
(\texttt{decn\_u\_complete}, then \texttt{decide5\_decides}).
\item \textbf{The CAD.} The run-time checks give delineability in the sign-invariant sense (\texttt{delin\_cl?}): if the
walk's certificate holds at one point of a connected set and the signs it reads there hold across the set, the family's
roots over the set are continuous and never cross, and the family has one sign vector on each section and band
(\texttt{stack\_const}). From this, the run's cells form a CAD (\texttt{cad\_run}), \texttt{cad\_out} returns them as
data that decides every sentence over \texttt{F} in the same variables (\texttt{cad\_out\_ok}, \texttt{cad\_decides}),
and every cell and root function's graph is first-order definable (\texttt{rcell\_fod\_le}, \texttt{graph\_fod});
\texttt{cad\_verified} puts them together.
\item \textbf{Quantifier elimination.} For every prenex formula with at least one quantifier and $m \ge 1$ free
variables, \texttt{qe} returns a quantifier-free formula that holds at exactly the same points of $\mathbb{R}^m$
(\texttt{qe\_complete}). With two or more free variables, when the plain run's answer is not quantifier-free, \texttt{qe}
falls back to a run whose free levels are closed under the derivative, where Thom's lemma separates the cells.
\texttt{(cad-qe)} uses \texttt{qe} in a theory that does not import \texttt{qe8}.
\item \textbf{Collins's projection.} Collins's N-reducta projection (\texttt{projn}, executable) is proved delineating
for every input, by the subresultant theorem and the continuity of complex roots in the coefficients
(\texttt{collins\_stack}, \texttt{projn\_ok}). The cells built from its tower form a CAD adapted to the family in every
number of variables, with no run-time certificate (\texttt{col\_cells}), which \texttt{col\_found} returns as data
(\texttt{col\_found\_cad}). \texttt{decide8} decides on that tower, and \texttt{qe8} eliminates quantifiers on it or on
one whose free levels are closed under the derivative (\texttt{decide8\_decides}, \texttt{qe8\_complete}).
\item \textbf{Tarski--Seidenberg.} \texttt{qelim} eliminates the quantifiers of every first-order formula from the inside
out, each existential quantifier by \texttt{qe8}, or by \texttt{decide8} when no variable is free (\texttt{qelim\_qf},
\texttt{qelim\_ok}). So a set is first-order definable over $\mathbb{Q}$ exactly when it is quantifier-free definable
(\texttt{fod\_qfd}), and every cell of both CADs is (\texttt{cad\_sa}).
\item \textbf{Formulas as written.} A datatype \texttt{Gm} mirrors PVS's connectives; \texttt{decide\_g} decides its
sentences and \texttt{qe\_g} eliminates its quantifiers, both through \texttt{qelim} (\texttt{decide\_g\_ok},
\texttt{qe\_g\_ok}). With \texttt{IMPORTING pvs\_cad}, the one import \texttt{(cad)} and \texttt{(cad-qe)} need, \texttt{(cad)}, \texttt{(cad *)}
and \texttt{(cad-qe)} take formulas as written: grouped binders, quantifiers inside NOT, AND, OR, IMPLIES and IFF, and
binders over \texttt{posreal}, \texttt{nnreal}, \texttt{negreal}, \texttt{npreal} and \texttt{nzreal} (a binder over any
other type, such as \texttt{nat}, is refused).
\end{enumerate}

\section{The library in numbers}
\begin{center}\small
\begin{tabular}{@{}l r@{\qquad}l r@{}}
\toprule
PVS theories & 355 & Formulas proved, whole-library replay & 5,183 of 5,183\\
Lines of PVS (358 files, with the datatypes) & 26,864 & \quad lemmas and theorems & 3,553\\
Lines of strategy code (Lisp) & 3,371 & \quad TCCs (type-correctness obligations) & 1,630\\
Commits, from 9 September 2026 & 298 & Axioms added & none\\
\bottomrule
\end{tabular}
\end{center}
It builds on NASALib's \texttt{Sturm} and \texttt{Tarski} (univariate root counting) and on \texttt{reals}, \texttt{structures}, \texttt{analysis}, \texttt{complex}, \texttt{mult\_poly}, \texttt{matrices} and \texttt{interval\_arith}.

\section{Where this sits}
Tarski showed that the first-order theory of the real numbers is decidable~\cite{tarski}, and Collins's
cylindrical algebraic decomposition~\cite{collins} made the decision practical; this library builds on the
PVS~\cite{pvs} libraries of NASALib~\cite{nasalib}, in particular the univariate Sturm and Tarski decision
procedures of Narkawicz, Mu\~noz and Dutle~\cite{nmd}. Earlier quantifier elimination for the reals with
machine-checked answers includes, by methods other than CAD, Cohen and Mahboubi's in Coq~\cite{cohenmahboubi},
McLaughlin and Harrison's proof-producing procedure in HOL Light~\cite{mclaughlin}, and a verified, complete
multivariate procedure in Isabelle/HOL~\cite{kosaian}, exported to executable code, whose authors report that without
virtual substitution it is impractical beyond the simplest univariate examples. Mahboubi implemented CAD in Coq without a
completed correctness proof~\cite{mahboubi}. A formal correctness proof of Collins's CAD, including the delineability
of Collins's projection, exists in Rocq~\cite{vermande}, over an abstract real closed field; it is a correctness proof,
not an executable procedure used as a tactic. Lean's Mathlib~\cite{mathlib} has, as far as we could find, no CAD
formalization and no CAD-based decision tactic, and a Lean project's roadmap from September 2026 lists CAD,
delineability and Sturm theory among what Mathlib lacks~\cite{hex}; Lean's nonlinear tactics (\texttt{nlinarith},
\texttt{polyrith}) are heuristic or check certificates found by an external algebra system, and are not complete. To
our knowledge (a survey of September 2026), pvs\_cad is the first multivariate decision procedure of the CAD family
that is executable, proof-producing, formally verified sound and complete, and usable as a tactic in an interactive
prover. The theorems about both its CADs (for the n-level engine, whenever its certificate holds) state the conditions
of Basu, Pollack and Roy's Def.~5.1, in two respects only indirectly: the partition is stated for the cells of
$\mathbb{R}^n$, and at lower levels it follows from the stack construction (\texttt{tcell\_cover}, \texttt{tcell\_disj})
without being restated; finiteness comes through the records of \texttt{cad\_out} and \texttt{col\_found}, at least one
for every nonempty cell. What is proved, and what is not:
\begin{itemize}[leftmargin=1.3em, itemsep=1pt, topsep=2pt]
\item ``Semi-algebraic'' holds in two senses. The CAD theorems state that every cell, at every level, and every root
function's graph over a cell is defined by a first-order formula whose atoms are sign conditions on polynomials with
rational coefficients. \texttt{fod\_qfd} (\texttt{qelim\_ok.pvs}) proves that a set is defined by such a formula exactly
when a Boolean combination of such sign conditions defines it without quantifiers (Tarski--Seidenberg over
$\mathbb{Q}$), and \texttt{cad\_sa.pvs} states this for every cell, at every level, of both CADs and for the Collins
CAD's root graphs (\texttt{col\_cells\_sa}, \texttt{cad\_cells\_sa}). Sets so defined over $\mathbb{Q}$ are
semi-algebraic in Basu, Pollack and Roy's sense, whose definition also allows real coefficients and so admits more sets
($\{\pi\}$, for one). No tool prints the quantifier-free definitions of the cells; the theorems say that they exist.
\item For the n-level engine, delineability is proved for the polynomial families the run builds and checks. For
Collins's N-reducta projection it is proved for every input (\texttt{collins\_stack}, \texttt{col\_stack.pvs});
\texttt{projn\_ok} ties it to the executable operator \texttt{projn}, whose tower \texttt{decide8} and \texttt{qe8} use.
McCallum's and Brown--McCallum's smaller projections are not proved delineating (the n-level engine starts its read
closure from a McCallum-shaped operator, \texttt{cad\_proj}, but its correctness does not rest on it), and Lazard's is
not done. Delineability is in the sign-invariant sense (\texttt{delin\_cl?}): root multiplicities are not tracked.
\item The Collins CAD theorems cover every number of variables, one included (\texttt{col\_cells}; in one variable the
cells are the sectors of the real roots of \texttt{F}, \texttt{col\_cad\_line}, the sectors over which
\texttt{decide8} and \texttt{decide5} decide one-variable sentences); the n-level engine's theorems (\texttt{cad\_verified}) are stated for two or more.
\item The CAD theorems cover the n-level engine and the Collins run. \texttt{(cad)} also has shortcuts they do not
cover: the witness search and, with \texttt{decide5}, the two-level method \texttt{decw} for two variables; their answers
are proved correct too.
\texttt{(cad :cad-only? t)} uses neither: it decides by \texttt{decide8} in a theory that imports \texttt{cad\_decide8}
and by \texttt{decide7} (the n-level engine) otherwise, and by \texttt{decide\_g} for formulas of other shapes.
\end{itemize}

\section{The main theorems}
A prenex sentence is encoded as data: \texttt{os} lists the quantifiers outermost first (\texttt{TRUE} for $\forall$,
\texttt{FALSE} for $\exists$), \texttt{F} lists the polynomials (exact rational coefficients), and \texttt{phi} is the
Boolean skeleton, whose atoms say ``polynomial $i$ has sign $s$''. Its meaning is defined by recursion on the prefix, with
the real quantifiers of PVS itself:

\begin{thmbox}{The semantics \Cm{\normalfont(cad\_decide.pvs)}}
\begin{Verbatim}
fsem(os, F, phi, ys): \Kw{RECURSIVE} bool =
  \Kw{CASES} os \Kw{OF}
    null: bfeval(phi, F, ys),
    cons(q, r): \Kw{IF} q \Kw{THEN} \Kw{FORALL} x: fsem(r, F, phi, cons(x, ys))
                \Kw{ELSE} \Kw{EXISTS} x: fsem(r, F, phi, cons(x, ys)) \Kw{ENDIF}
  \Kw{ENDCASES}
\end{Verbatim}
\end{thmbox}

The decision procedures return a record: \texttt{val}, the verdict, and \texttt{ok}, a flag that is TRUE when the
run-time checks passed. \texttt{decide8} is the Collins run; \texttt{decide5} is built on the n-level engine:

\begin{thmbox}{Soundness and completeness \Cm{\normalfont(cad\_decide8.pvs; cad\_decide5.pvs, complete\_all.pvs)}}
\begin{Verbatim}
decide8_correct:  \Kw{THEOREM} decide8(reverse(os), F, phi)`ok \Kw{IMPLIES}
                    (decide8(reverse(os), F, phi)`val \Kw{IFF} fsem(os, F, phi, null))
decide8_decides:  \Kw{THEOREM} decide8(reverse(os), F, phi)`val \Kw{IFF} fsem(os, F, phi, null)
decide5_correct:  \Kw{THEOREM} decide5(reverse(os), F, phi)`ok \Kw{IMPLIES}
                    (decide5(reverse(os), F, phi)`val \Kw{IFF} fsem(os, F, phi, null))
decide5_complete: \Kw{THEOREM} decide5(reverse(os), F, phi)`ok
decide5_decides:  \Kw{THEOREM} decide5(reverse(os), F, phi)`val \Kw{IFF} fsem(os, F, phi, null)
\end{Verbatim}
\end{thmbox}

There are no hypotheses: \texttt{os}, \texttt{F} and \texttt{phi} are arbitrary, so the theorems cover every closed
prenex sentence with any number of quantifiers. Both procedures are total functions, and PVS generated and proved their
termination obligations, so each always terminates and returns the truth of the sentence. \texttt{(cad :cad-only? t)}
uses \texttt{decide8} in a theory that imports \texttt{cad\_decide8}, otherwise \texttt{decide7}
(\texttt{cad\_decide7.pvs}), whose \texttt{decide7\_correct} and \texttt{decide7\_decides} have the statements above
with \texttt{decide7} in place of \texttt{decide5}.

Quantifiers anywhere need another formula type, \texttt{rcf\_fol}'s \texttt{Fm} (\texttt{ftrue}, sign atoms
\texttt{fsg(p, s)}, \texttt{fnot}, \texttt{fand}, \texttt{fex}; variables are positions in a point). Its semantics
\texttt{fsem(phi)(pt)} is a different function from \texttt{fsem(os, F, phi, ys)} above; \texttt{qfsem} is the same
semantics under another name, and \texttt{qf?(phi)} says that \texttt{phi} has no \texttt{fex}. \texttt{gform\_def}'s
\texttt{Gm} mirrors PVS's connectives (TRUE, FALSE, the six comparisons of two polynomials, NOT, AND, OR, IMPLIES, IFF, FORALL,
EXISTS) and its meaning \texttt{gsem(g)(pt)} is written with them; \texttt{decide\_g(g)} and \texttt{qe\_g(g, m)} apply
\texttt{qelim} to its translation into \texttt{Fm}. \texttt{qelim} works from the inside out: each \texttt{fex} by
\texttt{qe8}, or by \texttt{decide8} when no variable is free.

In the theorems below on \texttt{g} and \texttt{qelim} the formula is arbitrary and \texttt{m}, where it occurs, is any
natural number (\texttt{m = 0}: a sentence). In \texttt{qe8\_complete}, \texttt{m} is a positive natural number and \texttt{cons?(os)}
asks for at least one quantifier: every prenex formula with at least one quantifier and $m \ge 1$ free variables;
\texttt{qe\_complete} (\texttt{qe\_all.pvs}) states the same for \texttt{qe}, built on the n-level engine. The atoms'
polynomials have rational coefficients and the formulas have no parameters, so \texttt{fod\_qfd} is the
Tarski--Seidenberg theorem for formulas over $\mathbb{Q}$.

\begin{thmbox}{Any shape, quantifier elimination, Tarski--Seidenberg \Cm{\normalfont(gform\_ok, qe8, qelim\_ok, rcf\_fol)}}
\begin{Verbatim}[fontsize=\footnotesize]
decide_g_ok:  \Kw{THEOREM} decide_g(g) \Kw{IFF} gsem(g)(null)
qe_g_qf:      \Kw{LEMMA} qf?(qe_g(g, m))
qe_g_ok:      \Kw{THEOREM} length(pt) = m \Kw{IMPLIES} (fsem(qe_g(g, m))(pt) \Kw{IFF} gsem(g)(pt))
qe8_complete: \Kw{THEOREM} cons?(os) \Kw{AND} length(pt) = m \Kw{IMPLIES}
    qf?(qe8(m, os, F, phi)`out) \Kw{AND} (qfsem(qe8(m, os, F, phi)`out)(pt) \Kw{IFF} fsem(os, F, phi, pt))
qelim_qf:     \Kw{LEMMA} \Kw{FORALL} phi: \Kw{FORALL} m: qf?(qelim(phi, m))
qelim_ok:     \Kw{THEOREM} \Kw{FORALL} phi: \Kw{FORALL} m, pt: length(pt) = m
    \Kw{IMPLIES} (fsem(qelim(phi, m))(pt) \Kw{IFF} fsem(phi)(pt))
fod?(n)(S): bool = \Kw{EXISTS} phi: \Kw{FORALL} pt: length(pt) = n \Kw{IMPLIES} (S(pt) \Kw{IFF} fsem(phi)(pt))
qfd?(n)(S): bool = \Kw{EXISTS} phi: qf?(phi) \Kw{AND} \Kw{FORALL} pt: length(pt) = n \Kw{IMPLIES} (S(pt) \Kw{IFF} fsem(phi)(pt))
fod_qfd:      \Kw{THEOREM} fod?(n)(S) \Kw{IFF} qfd?(n)(S)
\end{Verbatim}
\end{thmbox}

\begin{thmbox}{The Collins CAD \Cm{\normalfont(col\_stack.pvs, col\_line.pvs, col\_found\_ok.pvs, cad\_sa.pvs)}}
\begin{Verbatim}[fontsize=\footnotesize]
collins_stack: \Kw{THEOREM} conn?(C) \Kw{AND} C(p0) \Kw{AND} colh?(G, C) \Kw{IMPLIES} delin_cl?(G, C)
col_cells:     \Kw{THEOREM} ccells?(k, F)
col_found_cad: \Kw{THEOREM} cons?(qs) \Kw{IMPLIES} cfound?(qs, F, col_found(qs, F)`frecs)
col_cells_sa: \Kw{THEOREM}
      (\Kw{FORALL} s, A: member(s, csects(k, F)) \Kw{AND} length(A) <= k
         \Kw{IMPLIES} qfd?(1 + length(A))(ccell(k, F)(s, A)))
  \Kw{AND} (\Kw{FORALL} s, A, i: member(s, csects(k, F)) \Kw{AND} length(A) < k
         \Kw{AND} (\Kw{FORALL} q: ccell(k, F)(s, A)(q) \Kw{IMPLIES} i < nroot(nth(ctw(F, k), length(A)), q))
         \Kw{IMPLIES} qfd?(2 + length(A))
           (\Kw{LAMBDA} p: cons?(p) \Kw{AND} ccell(k, F)(s, A)(cdr(p))
                      \Kw{AND} car(p) = rt(nth(ctw(F, k), length(A)), cdr(p), i)))
\end{Verbatim}
\end{thmbox}

\texttt{colh?(G, C)} says that Collins's projection of \texttt{G} has one sign on \texttt{C}: the leading coefficients
of the kept reducta, their principal subresultant coefficients with their derivatives, and those of pairs of members;
\texttt{projn\_ok} (\texttt{cad\_projn.pvs}) derives it from the sign-invariance of the executable operator
\texttt{projn}. \texttt{ccells?(k, F)} states, for the cells \texttt{ccell(k, F)(s, A)} built from the tower
\texttt{ctw(F, k)} of Collins projections of a family \texttt{F} in $k+1$ variables, that the cells of $\mathbb{R}^{k+1}$
cover it and are pairwise disjoint, that every cell at every level is connected and first-order definable, that over
each sector of the line every level is delineable over every cell below it, that the graph of every root function over a
cell is definable, and that \texttt{F} has one sign vector on every cell of $\mathbb{R}^{k+1}$; \texttt{col\_cells} has no
hypothesis: the cells form a CAD with no run-time check (\texttt{col\_found} and \texttt{decide8} still check their walk's
certificate while they raise $u$, which always ends, \texttt{decc\_complete}). \texttt{cfound?(qs, F, L)} adds the records \texttt{L}: each lies in its cell, with an
exact sample point and \texttt{F}'s signs there, and every nonempty cell of $\mathbb{R}^n$, $n$ = \texttt{length(qs)},
has at least one; by \texttt{col\_found\_decides} they decide every sentence over \texttt{F} in the same variables, in
the same order. \texttt{col\_cells\_sa} states that every cell, at every level, and every root function's graph over a
cell is quantifier-free definable.

\begin{thmbox}{The n-level engine's CAD \Cm{\normalfont(cad\_verified.pvs, cad\_sa.pvs)}}
\begin{Verbatim}[fontsize=\footnotesize]
cad_verified: \Kw{THEOREM} decn_o(u, qs, F, Psi)`ok \Kw{AND} length(qs) >= 2 \Kw{IMPLIES} cad_of?(u, qs, F)
cad_exists:   \Kw{THEOREM} length(qs) >= 2 \Kw{IMPLIES} \Kw{EXISTS} u: cad_of?(u, qs, F)
cad_cells_sa: \Kw{THEOREM} cad_of?(u, qs, F) \Kw{AND} member(s, runsects(u, length(cdr(qs)), F))
  \Kw{AND} length(A) <= length(cdr(qs))
  \Kw{IMPLIES} qfd?(1 + length(A))(rcell(u, length(cdr(qs)), F)(s, A))
\end{Verbatim}
\end{thmbox}

\texttt{cad\_of?(u, qs, F)} states the same conditions as \texttt{ccells?} for the cells \texttt{rcell(u, k, F)(s, A)}
of the n-level engine \texttt{decn\_o}, where $k+1$ = \texttt{length(qs)} and each sector of the line has a tower of its
own, with the records of \texttt{cad\_out}: at least one, with a sample point and
\texttt{F}'s signs, for every nonempty cell of $\mathbb{R}^n$, and none for any other. \texttt{cad\_cells\_sa}: whenever
\texttt{cad\_of?} holds, every cell, at every level, is quantifier-free definable.


\begin{thebibliography}{99}\footnotesize
\setlength\itemsep{1pt}
\bibitem{tarski} A.~Tarski. \emph{A Decision Method for Elementary Algebra and Geometry}. RAND Corporation, 1948;
  2nd ed., University of California Press, 1951.
\bibitem{collins} G.~E. Collins. Quantifier elimination for real closed fields by cylindrical algebraic
  decomposition. In \emph{Automata Theory and Formal Languages}, LNCS 33, pages 134--183. Springer, 1975.
\bibitem{pvs} S.~Owre, J.~M. Rushby, N.~Shankar. PVS: A prototype verification system. In \emph{CADE-11},
  LNAI 607. Springer, 1992. \url{https://pvs.csl.sri.com}
\bibitem{nasalib} NASALib, the NASA PVS library collection. \url{https://github.com/nasa/pvslib}
\bibitem{nmd} A.~Narkawicz, C.~Mu\~noz, A.~Dutle. Formally-verified decision procedures for univariate
  polynomial computation based on Sturm's and Tarski's theorems. \emph{Journal of Automated Reasoning}
  54:285--326, 2015. \url{https://doi.org/10.1007/s10817-015-9320-x}
\bibitem{cohenmahboubi} C.~Cohen, A.~Mahboubi. Formal proofs in real algebraic geometry: from ordered fields to
  quantifier elimination. \emph{Logical Methods in Computer Science} 8(1:02), 2012.
\bibitem{mclaughlin} S.~McLaughlin, J.~Harrison. A proof-producing decision procedure for real arithmetic. In
  \emph{CADE-20}, LNCS 3632, pages 295--314. Springer, 2005.
\bibitem{kosaian} K.~Kosaian, Y.~K. Tan, A.~Platzer. A first complete algorithm for real quantifier
  elimination in Isabelle/HOL. In \emph{CPP 2023}. ACM, 2023. \url{https://arxiv.org/abs/2209.10978}
\bibitem{mahboubi} A.~Mahboubi. Implementing the cylindrical algebraic decomposition within the Coq system.
  \emph{Mathematical Structures in Computer Science} 17(1), 2007.
\bibitem{vermande} Q.~Vermande. Cylindrical algebraic decomposition in Coq/Rocq. In \emph{CPP 2026},
  pages 45--58. ACM, 2026. \url{https://doi.org/10.1145/3779031.3779100}
\bibitem{mathlib} The mathlib Community. The Lean mathematical library. In \emph{CPP 2020}. ACM, 2020.
\bibitem{hex} Hex project (Lean), roadmap for the real algebraic geometry behind CAD, issue 10300,
  September 2026. \url{https://github.com/kim-em/hex-dev/issues/10300}
\bibitem{motzkin} T.~S. Motzkin. The arithmetic-geometric inequality. In O.~Shisha (ed.), \emph{Inequalities},
  pages 205--224. Academic Press, 1967.
\end{thebibliography}

\end{document}
